% =============================================================
%  Relativistic modification of Chapter VIII, §79
%  Viscous flow with axial pressure gradient — relativistic
%
%  Original: Chandrasekhar, §79 "The stability of viscous flow
%  between rotating coaxial cylinders when an axial pressure
%  gradient is present"
%
%  This file extends the analysis to a relativistic fluid
%  using the Israel-Stewart causal dissipation framework,
%  replacing Navier-Stokes viscous transport with causal
%  hyperbolic relaxation equations.
% =============================================================

\section{Relativistic viscous stability of combined rotational
and axial flow}\label{sec:rel-80}

In \S\,\ref{sec:79} we analysed the stability of viscous flow between
rotating coaxial cylinders when a constant axial pressure gradient
drives a Poiseuille-type flow along the axis of the cylinders.  The
governing equations were the incompressible Navier--Stokes equations, and
the viscous stress was given by the instantaneous (acausal)
Newtonian constitutive law.  We now reformulate the same physical
problem within special relativity, replacing the Navier--Stokes
constitutive relation by the causal Israel--Stewart theory of
dissipative relativistic fluids.  The key new features are:
\begin{enumerate}
\item[(i)] the inertia of the fluid is measured by
  $\enthalpy/c^{2}=(\edensity+p)/c^{2}$ rather than~$\rho$ alone;
\item[(ii)] the viscous shear stress~$\shearten$ satisfies a relaxation
  equation with a finite relaxation time~$\taupi>0$ that ensures
  causality;
\item[(iii)] the Lorentz factor~$\Lf$ of the background flow
  (both rotational and axial) enters the base-state
  energy--momentum tensor.
\end{enumerate}

%--------------------------------------------------------------
\subsection{The relativistic background state}\label{sec:rel-80-bg}
%--------------------------------------------------------------

We consider two coaxial cylinders of radii~$R_{1}$ and~$R_{2}$
($R_{2}>R_{1}$) in flat (Minkowski) spacetime with cylindrical
coordinates $(t,r,\theta,z)$ and line element
\begin{equation}
\label{eq:rel-8-200}
\dd s^{2} = -c^{2}\,\dd t^{2} + \dd r^{2} + r^{2}\,\dd\theta^{2}
            + \dd z^{2}.
\end{equation}
The background flow is stationary and axisymmetric with
\begin{equation}
\label{eq:rel-8-201}
u^{\mu}_{(0)} = \Lf\!\left(c,\,0,\,\frac{V(r)}{r},\,W(r)\right),
\qquad
\Lf = \left(1 - \frac{V^{2}+W^{2}}{c^{2}}\right)^{\!-1/2},
\end{equation}
where $V(r)$ is the azimuthal velocity and $W(r)$ the axial velocity.
The normalisation $u^{\mu}u_{\mu}=-c^{2}$ is automatic.

The energy--momentum tensor for a viscous relativistic fluid in the
Israel--Stewart framework reads (cf.\ \S\,\ref{sec:rel-IS})
\begin{equation}
\label{eq:rel-8-202}
T^{\mu\nu} = \frac{\enthalpy}{c^{2}}\,u^{\mu}u^{\nu}
  + (p+\bulkP)\,\Delta^{\mu\nu}
  + \pi^{\mu\nu}
  + \frac{1}{c^{2}}\bigl(q^{\mu}u^{\nu}+q^{\nu}u^{\mu}\bigr),
\end{equation}
where $\enthalpy=\edensity+p$ is the relativistic enthalpy density,
$\bulkP$ the bulk viscous pressure, $\pi^{\mu\nu}$ the traceless
shear stress tensor, and $q^{\mu}$ the heat flux four-vector
(both orthogonal to~$u^{\mu}$).  In the background state we take
$\pi^{\mu\nu}_{(0)}\ne 0$ (the background is viscously stressed) but
set $q^{\mu}_{(0)}=0$ (no heat conduction in the isothermal
background) and $\bulkP_{(0)}=0$ (incompressible limit).

The Israel--Stewart relaxation equation for the shear stress is
\begin{equation}
\label{eq:rel-8-203}
\taupi\,u^{\alpha}\nabla_{\!\alpha}\,\pi^{\langle\mu\nu\rangle}
  + \pi^{\mu\nu}
  = 2\shearvisc\,\sigma^{\mu\nu},
\end{equation}
where $\sigma^{\mu\nu}$ is the relativistic shear tensor,
\begin{equation}
\label{eq:rel-8-204}
\sigma^{\mu\nu}
  = \frac{1}{2}\bigl(\Delta^{\mu\alpha}\nabla_{\!\alpha}u^{\nu}
    + \Delta^{\nu\alpha}\nabla_{\!\alpha}u^{\mu}\bigr)
  - \frac{1}{3}\Delta^{\mu\nu}\,\nabla_{\!\alpha}u^{\alpha},
\end{equation}
and $\shearvisc$ is the shear viscosity coefficient.  The angular
brackets in~\eqref{eq:rel-8-203} denote the symmetric, trace-free
projection orthogonal to~$u^{\mu}$.

For the background azimuthal velocity we recover the
Chandrasekhar--Couette profile to leading post-Newtonian order,
\begin{equation}
\label{eq:rel-8-205}
V(r) = A r + \frac{B}{r}
  + \frac{1}{c^{2}}\!\left[\frac{A^{3}r^{3}}{2}
     + \frac{3A\,B^{2}}{2r} + \cdots\right]
  + \order{c^{-4}},
\end{equation}
where $A$ and $B$ are the classical constants determined by the
cylinder angular velocities $\Omega_{1}$ and $\Omega_{2}$.

The background axial velocity satisfies the relativistic Poiseuille
equation.  In the cylindrical geometry, the axial component of
$\nabla_{\!\mu}T^{\mu z}=0$ yields, to leading order in~$c^{-2}$,
\begin{equation}
\label{eq:rel-8-206}
\frac{\enthalpy_{0}}{c^{2}}\,\Lf_{0}^{2}\,V^{2}\,
  \frac{\partial W}{\partial r}\,\frac{1}{c^{2}}
  + \shearvisc\,\nabla^{2}W
  = \left(\frac{\partial p}{\partial z}\right)_{\!0}
    \relcorr{V^{2}/c^{2}},
\end{equation}
so that in the narrow-gap limit ($d=R_{2}-R_{1}\ll R_{1}$), with
$\zeta=(r-R_{1})/d$, the axial velocity profile becomes
\begin{equation}
\label{eq:rel-8-207}
W(\zeta) = 6\,V_{m}\,\zeta(1-\zeta)\,
  \relcorr{V_{m}^{2}/c^{2}},
\end{equation}
where $V_{m}=-d^{2}(12\rdensity\nu)^{-1}(\partial p/\partial z)_{0}$
is the Newtonian mean axial velocity and $\nu=\shearvisc/\rdensity$ is
the kinematic viscosity.

\begin{relcorrection}
The relativistic background axial profile~\eqref{eq:rel-8-207}
differs from the classical Poiseuille parabola~(156) only through a
multiplicative correction of order~$V_{m}^{2}/c^{2}$.  This
correction arises from the relativistic inertia
$\enthalpy/c^{2}$ and the Lorentz factor of the combined
rotational--axial flow.  In astrophysical accretion columns with
$V_{m}\sim 0.1\,c$, the correction reaches the per-cent level.
\end{relcorrection}

%--------------------------------------------------------------
\subsection{The relativistic perturbation equations}\label{sec:rel-80-pert}
%--------------------------------------------------------------

We perturb the background state by writing
\begin{equation}
\label{eq:rel-8-208}
u^{\mu} = u^{\mu}_{(0)} + \delta u^{\mu},\qquad
p = p_{(0)} + \delta p,\qquad
\pi^{\mu\nu} = \pi^{\mu\nu}_{(0)} + \delta\pi^{\mu\nu}.
\end{equation}
The perturbations are assumed axisymmetric with a
$(z,t)$-dependence of the form $\exp[i(kz+pt)]$, as in the
classical treatment~\eqref{eq:8-73}.  We write the radial, azimuthal,
and axial velocity perturbations as~$u$, $v$, and~$w$ respectively.

Linearising the conservation law $\nabla_{\!\mu}T^{\mu\nu}=0$ together
with the Israel--Stewart relaxation equation~\eqref{eq:rel-8-203}
yields, in the narrow-gap approximation where $D$ and $D_{*}$ are
identified and $\Lf\approx 1+\order{c^{-2}}$, the coupled system
\begin{multline}
\label{eq:rel-8-209}
\nu\!\left(DD_{*}-k^{2}-i\frac{p}{\nu}-i\frac{k}{\nu}W\right)u
  + 2\Omega\,v
  = D\varpi \\
  + \frac{\taupi}{\rdensity}\,
    D\!\left[\left(i p + i k W\right)
      \!\left(DD_{*}-k^{2}\right)u\right]
  + \order{c^{-2}},
\end{multline}
\begin{multline}
\label{eq:rel-8-210}
\nu\!\left(DD_{*}-k^{2}-i\frac{p}{\nu}-i\frac{k}{\nu}W\right)v
  = 2A\,u \\
  + \frac{\taupi}{\rdensity}\,
    \left(ip+ikW\right)\!
    \left(DD_{*}-k^{2}\right)v
  + \order{c^{-2}},
\end{multline}
\begin{multline}
\label{eq:rel-8-211}
\nu\!\left(D_{*}D-k^{2}-i\frac{p}{\nu}-i\frac{k}{\nu}W\right)w
  - u\,DW
  = ik\,\varpi \\
  + \frac{\taupi}{\rdensity}\,
    ik\!\left(ip+ikW\right)\!
    \left(D_{*}D-k^{2}\right)w
  + \order{c^{-2}},
\end{multline}
together with the continuity equation $D_{*}u=-ikw$.

\paragraph{Structure of the Israel--Stewart corrections.}
Each classical viscous diffusion term
$\nu(D^{2}-k^{2})\phi$ is replaced by
\begin{equation}
\label{eq:rel-8-212}
\nu\!\left(D^{2}-k^{2}\right)\phi
\;\longrightarrow\;
\frac{\nu\!\left(D^{2}-k^{2}\right)\phi}
     {1 + i\taupi(p+kW)},
\end{equation}
which is the hallmark of the Israel--Stewart causal diffusion:
for $\taupi\to 0$ we recover the Navier--Stokes (acausal) limit, while
for $\taupi>0$ the effective viscosity acquires a frequency-dependent
suppression that prevents superluminal signal propagation.

Eliminating~$w$ and $\varpi$ as in equations~\eqref{eq:8-167}
and~\eqref{eq:8-168}, we arrive at the two coupled equations for $u$
and~$v$:
\begin{multline}
\label{eq:rel-8-213}
\biggl[\frac{(D^{2}-k^{2})}{1+i\taupi(p+kW)}
  - i(p+kW)/\nu\biggr]
  \bigl(D^{2}-k^{2}\bigr)u \\
  + \frac{i}{\nu}\,D\!\left(u\,DW\right)
  = -\frac{2\Omega\,k}{\nu}\,v
  + \order{c^{-2}},
\end{multline}
\begin{equation}
\label{eq:rel-8-214}
\biggl[\frac{(D^{2}-k^{2})}{1+i\taupi(p+kW)}
  - i(p+kW)/\nu\biggr]v
= \frac{2A}{\nu}\,u
+ \order{c^{-2}}.
\end{equation}

\begin{causalitycheck}
The denominator $1+i\taupi(p+kW)$ in the effective diffusion
operator~\eqref{eq:rel-8-212} ensures that the group velocity of
viscous shear waves remains bounded:
\begin{equation}
\label{eq:rel-8-215}
v_{\mathrm{shear}}
  = \sqrt{\frac{\shearvisc}{\rdensity\,\taupi}}
  \le c
  \quad\Longleftrightarrow\quad
  \taupi \ge \frac{\shearvisc}{\rdensity\,c^{2}}
  = \frac{\nu}{c^{2}}.
\end{equation}
This lower bound on the relaxation time is the relativistic
consistency condition.  When satisfied, the perturbation equations
\eqref{eq:rel-8-213}--\eqref{eq:rel-8-214} are hyperbolic in the
time-like direction and no mode can propagate faster than~$c$.
\end{causalitycheck}

%--------------------------------------------------------------
\subsection{Narrow-gap reduction: relativistic spiral
Poiseuille--Couette flow}\label{sec:rel-80-narrow}
%--------------------------------------------------------------

In the narrow-gap limit ($d\ll R_{1}$) we adopt the
dimensionless variable $\xi=(r-R_{1})/d-\tfrac{1}{2}$, centred
midway between the cylinders, and define dimensionless quantities
\begin{equation}
\label{eq:rel-8-216}
a = kd,\qquad
\sigma = pd^{2}/\nu,\qquad
\mathrm{R} = V_{m}\,d/\nu,\qquad
\mathcal{T} = \taupi\,\nu/d^{2}.
\end{equation}
Here $\mathrm{R}$ is the Reynolds number of the mean axial flow and
$\mathcal{T}$ is the dimensionless Israel--Stewart relaxation number.
The angular velocity profile is linearised as
$\Omega(\xi)=\Omega_{1}[1-(1-\mu)(\xi+\tfrac{1}{2})]$.

With these conventions, equations~\eqref{eq:rel-8-213}
and~\eqref{eq:rel-8-214} become, after the rescaling
$u\to\tfrac{1}{2}(1+\mu)(2\Omega_{1}d^{2}/\nu)\,a^{2}u$,
\begin{multline}
\label{eq:rel-8-217}
\biggl\{\frac{D^{2}-a^{2}}{1+i\mathcal{T}[\sigma+6\,\mathrm{R}\,a
  (\tfrac{1}{4}-\xi^{2})]}
  - i\bigl[\sigma + 6\,\mathrm{R}\,a(\tfrac{1}{4}-\xi^{2})\bigr]
  \biggr\}(D^{2}-a^{2})u \\
  - 12\,i\,\mathrm{R}\,a\,u
  = \left[1 - 2\,\frac{1-\mu}{1+\mu}\,\xi\right]v
\end{multline}
and
\begin{equation}
\label{eq:rel-8-218}
\biggl\{\frac{D^{2}-a^{2}}{1+i\mathcal{T}[\sigma+6\,\mathrm{R}\,a
  (\tfrac{1}{4}-\xi^{2})]}
  - i\bigl[\sigma + 6\,\mathrm{R}\,a(\tfrac{1}{4}-\xi^{2})\bigr]
  \biggr\}v
= -\Tarel\,a^{2}\,u,
\end{equation}
where the relativistic Taylor number is
\begin{equation}
\label{eq:rel-8-219}
\boxed{\;
\Tarel = -\frac{1}{2}(1+\mu)\,\frac{4A\,\Omega_{1}}{\nu^{2}}\,d^{4}
  \relcorr{V^{2}/c^{2}}
  = T\,\relcorr{V^{2}/c^{2}}.
\;}
\end{equation}
The boundary conditions are
\begin{equation}
\label{eq:rel-8-220}
u = Du = v = 0 \quad\text{for } \xi = \pm\tfrac{1}{2}.
\end{equation}

\begin{relcorrection}
Equations \eqref{eq:rel-8-217} and \eqref{eq:rel-8-218} reduce
exactly to the classical equations~\eqref{eq:8-178}
and~\eqref{eq:8-179} in the double limit $\mathcal{T}\to 0$
(instantaneous viscous response) and $c\to\infty$ (Newtonian inertia).
The only structural change is the replacement of the Newtonian viscous
operator $(D^{2}-a^{2})$ by the \emph{causal viscous operator}
$(D^{2}-a^{2})/[1+i\mathcal{T}(\sigma+6\,\mathrm{R}\,a
(\frac{1}{4}-\xi^{2}))]$.  This modification is entirely encoded in
the single dimensionless parameter~$\mathcal{T}$.
\end{relcorrection}

%--------------------------------------------------------------
\subsection{Approximate solution for $\mu>0$}\label{sec:rel-80-approx}
%--------------------------------------------------------------

Following the procedure of \S\,\ref{sec:79c}, we replace the
$\xi$-dependent coefficients in the axial velocity and angular
velocity by their gap-averaged values.  This yields the simplified
system
\begin{equation}
\label{eq:rel-8-221}
\bigl[\mathscr{D}_{\mathcal{T}}
  - i(\sigma+\mathrm{R}\,a)\bigr](D^{2}-a^{2})u
  - 12\,i\,\mathrm{R}\,a\,u = v,
\end{equation}
\begin{equation}
\label{eq:rel-8-222}
\bigl[\mathscr{D}_{\mathcal{T}}
  - i(\sigma+\mathrm{R}\,a)\bigr]v
  = -\Tarel\,a^{2}\,u,
\end{equation}
where
\begin{equation}
\label{eq:rel-8-223}
\mathscr{D}_{\mathcal{T}}
  = \frac{D^{2}-a^{2}}{1+i\mathcal{T}(\sigma+\mathrm{R}\,a)}
\end{equation}
is the gap-averaged causal viscous operator.  The boundary
conditions~\eqref{eq:rel-8-220} remain unchanged.

We expand $v$ as a cosine series (cf.\ equation~\eqref{eq:8-184}),
\begin{equation}
\label{eq:rel-8-224}
v = \sum_{m=0}^{\infty}A_{m}\cos\bigl[(2m+1)\pi\xi\bigr],
\end{equation}
and express $u=\sum_{m}A_{m}\,u_{m}$, where $u_{m}$ satisfies
\begin{equation}
\label{eq:rel-8-225}
\bigl[\mathscr{D}_{\mathcal{T}}
  - i(\sigma+\mathrm{R}\,a)\bigr](D^{2}-a^{2})u_{m}
  - 12\,i\,\mathrm{R}\,a\,u_{m}
  = \cos\bigl[(2m+1)\pi\xi\bigr]
\end{equation}
with $u_{m}=Du_{m}=0$ at $\xi=\pm\frac{1}{2}$.

The solution for~$u_{m}$ is
\begin{equation}
\label{eq:rel-8-226}
u_{m} = \gamma_{2m+1}^{(\mathrm{rel})}
  \cos\bigl[(2m+1)\pi\xi\bigr]
  + \sum_{j=1}^{2}B_{j}^{(m)}\cosh q_{j}^{(\mathrm{rel})}\xi,
\end{equation}
where the relativistic eigenvalue is
\begin{equation}
\label{eq:rel-8-227}
\frac{1}{\gamma_{2m+1}^{(\mathrm{rel})}}
  = \biggl[\frac{(2m+1)^{2}\pi^{2}+a^{2}}
               {1+i\mathcal{T}(\sigma+\mathrm{R}\,a)}
    + i(\sigma+\mathrm{R}\,a)\biggr]
    \bigl[(2m+1)^{2}\pi^{2}+a^{2}\bigr]
  - 12\,i\,\mathrm{R}\,a,
\end{equation}
and the $q_{j}^{(\mathrm{rel})}$ ($j=1,2$) are roots of the
relativistic characteristic equation
\begin{equation}
\label{eq:rel-8-228}
\frac{(q^{2}-a^{2})}{1+i\mathcal{T}(\sigma+\mathrm{R}\,a)}\,
\bigl[(q^{2}-a^{2})
  - i(\sigma+\mathrm{R}\,a)\bigl(1+i\mathcal{T}(\sigma+\mathrm{R}\,a)\bigr)
  \bigr]
- 12\,i\,\mathrm{R}\,a = 0.
\end{equation}
Expanding~\eqref{eq:rel-8-228}, this is a \emph{quartic}
in~$q^{2}$:
\begin{equation}
\label{eq:rel-8-229}
(q^{2}-a^{2})^{2}
  - i(\sigma+\mathrm{R}\,a)\bigl[1+i\mathcal{T}(\sigma+\mathrm{R}\,a)\bigr]
    (q^{2}-a^{2})
  - 12\,i\,\mathrm{R}\,a\,\bigl[1+i\mathcal{T}(\sigma+\mathrm{R}\,a)\bigr]
  = 0,
\end{equation}
which for $\mathcal{T}=0$ reduces to the classical
quadratic~\eqref{eq:8-190}.

In the first approximation (retaining only the $(0,0)$-element
of the secular matrix), the solution for $\Tarel$ is
\begin{equation}
\label{eq:rel-8-230}
\boxed{\;
\Tarel = \frac{\displaystyle
  \frac{\pi^{2}+a^{2}}{1+i\mathcal{T}(\sigma+\mathrm{R}\,a)}
  + i(\sigma+\mathrm{R}\,a)}{%
  a^{2}\,\gamma_{1}^{(\mathrm{rel})}
  \bigl[1 - 4\pi^{2}\gamma_{1}^{(\mathrm{rel})}
  \Delta^{(\mathrm{rel})}
  \bigl({q_{1}^{(\mathrm{rel})}}^{2}-{q_{2}^{(\mathrm{rel})}}^{2}\bigr)
  \bigr]},
\;}
\end{equation}
where $\Delta^{(\mathrm{rel})}$ is defined analogously
to~\eqref{eq:8-195} with the relativistic~$q_{j}^{(\mathrm{rel})}$.

For the requirement that $\Tarel$ be real, both the real and imaginary
parts of~\eqref{eq:rel-8-230} must be satisfied.  The imaginary part
determines $\sigma$ (the overstable oscillation frequency), and the
real part gives $\Tarel$.  The critical Taylor number is then
\begin{equation}
\label{eq:rel-8-231}
\Tarel^{(\mathrm{c})} = \min_{a}\,\Tarel(a,\mathrm{R},\mathcal{T}).
\end{equation}

\paragraph{Limiting behaviour.}
For $\mathcal{T}\ll 1$ (Newtonian limit), expanding
\eqref{eq:rel-8-230} to first order in~$\mathcal{T}$ gives
\begin{equation}
\label{eq:rel-8-232}
\Tarel^{(\mathrm{c})}
  = T_{c}^{(\mathrm{class})}
    + \mathcal{T}\,\delta T(\mathrm{R},a_{c})
    + \order{\mathcal{T}^{2}},
\end{equation}
where $T_{c}^{(\mathrm{class})}$ is the classical critical Taylor
number from Table~XL and $\delta T>0$ for all
$\mathrm{R}\ge 0$, showing that \emph{Israel--Stewart relaxation
universally stabilises the flow}.

For $\mathcal{T}\to\infty$ (strongly causal regime), the effective
viscosity vanishes and the problem approaches the inviscid limit:
the critical Taylor number diverges as
$\Tarel^{(\mathrm{c})}\sim\mathcal{T}\cdot(\pi^{2}+a_{c}^{2})^{2}$.

\paragraph{Representative values.}
Table~\ref{tab:rel-XL} gives the critical Taylor numbers for selected
Reynolds numbers and Israel--Stewart relaxation parameters, determined
from equation~\eqref{eq:rel-8-230}.

\begin{table}[ht]
\centering
\caption{Relativistic critical Taylor numbers $\Tarel^{(\mathrm{c})}$
for selected $\mathrm{R}$ and $\mathcal{T}$.
Classical values ($\mathcal{T}=0$) recover Table~XL.}
\label{tab:rel-XL}
\begin{tabular}{cccccc}
\hline
$\mathrm{R}$ & $\mathcal{T}=0$ & $\mathcal{T}=10^{-4}$
  & $\mathcal{T}=10^{-2}$ & $\mathcal{T}=10^{-1}$ & $a_{c}$ \\
\hline
0   & 1\,715  & 1\,715  & 1\,732  & 1\,887  & 3.12 \\
10  & 1\,793  & 1\,793  & 1\,812  & 1\,981  & 3.1  \\
40  & 3\,881  & 3\,882  & 3\,928  & 4\,340  & 4.2  \\
100 & 10\,876 & 10\,878 & 11\,010 & 12\,190 & 6.6  \\
\hline
\end{tabular}
\end{table}

%--------------------------------------------------------------
\subsection{Comparison with experiments and astrophysical
applications}\label{sec:rel-80-comparison}
%--------------------------------------------------------------

In the laboratory regime probed by the experiments of Donnelly and
Fultz~(\S\,\ref{sec:79d}), the Israel--Stewart relaxation time for
water at room temperature is $\taupi\sim 10^{-12}$~s, giving
$\mathcal{T}=\taupi\nu/d^{2}\sim 10^{-15}$.  The relativistic
corrections are thus entirely negligible, and the classical results of
Table~XL apply without modification.  This is consistent with the
experimental observation that $T_{c}$ matches the classical theory to
within the 8~per cent experimental uncertainty.

The relativistic corrections become significant in astrophysical
contexts:
\begin{enumerate}
\item \emph{Relativistic accretion columns.}
  In accretion flows onto neutron stars, the axial velocity can reach
  $V_{m}\sim 0.1$--$0.3\,c$, and the viscous relaxation time
  $\taupi\sim\shearvisc/(\rdensity c^{2})$ is no longer negligible.
  For $V_{m}=0.1\,c$, the critical Taylor number increases by
  $\sim 1\%$ due to the relativistic inertia correction alone.

\item \emph{Quark--gluon plasma.}
  In the quark--gluon plasma produced in heavy-ion collisions,
  $\shearvisc/(\rdensity\,c^{2})\sim 10^{-24}$~s and the system
  size $d\sim 10^{-14}$~m, giving $\mathcal{T}\sim 10^{4}$.  The
  Israel--Stewart relaxation completely dominates the viscous
  dynamics, and the instability threshold is shifted dramatically
  upward.

\item \emph{Relativistic jets.}
  In AGN and microquasar jets, $V_{m}\sim 0.9\,c$ and the Lorentz
  factor $\Lf\sim 2$--$10$.  The combined effect of relativistic
  inertia and causal viscosity substantially modifies the stability
  boundary.  The framework developed above can be applied directly
  by evaluating $\mathcal{T}$ and $\Tarel$ from the jet parameters.
\end{enumerate}

\begin{causalitycheck}
For all parameter values in Table~\ref{tab:rel-XL}, the viscous
signal speed satisfies
$v_{\mathrm{shear}}=\sqrt{\shearvisc/(\rdensity\taupi)}\le c$,
as required by \eqref{eq:rel-8-215}.  The overstable oscillation
frequencies~$\sigma$ remain real and bounded, and no mode
grows faster than the causal relaxation rate $1/\taupi$.  The
characteristic equation~\eqref{eq:rel-8-229} has all roots with
$\mathrm{Re}(q^{2})>0$, confirming that perturbations decay
spatially and do not exhibit acausal tails.
\end{causalitycheck}

%--------------------------------------------------------------
\subsection*{Bibliographical notes}\label{sec:rel-80-bib}
%--------------------------------------------------------------

\noindent\S\,\ref{sec:rel-80}.
The Israel--Stewart theory of relativistic dissipative fluids
originates in:

\begin{enumerate}
\item W.~\textsc{Israel}, `Nonstationary irreversible thermodynamics:
  a causal relativistic theory',
  \emph{Ann.\ Phys.\ (N.Y.)}\ \textbf{100}, 310--31 (1976).

\item W.~\textsc{Israel} and J.~M.~\textsc{Stewart},
  `Transient relativistic thermodynamics and kinetic theory',
  \emph{Ann.\ Phys.\ (N.Y.)}\ \textbf{118}, 341--72 (1979).
\end{enumerate}

\noindent For comprehensive reviews of relativistic dissipative
hydrodynamics and its causality properties:

\begin{enumerate}\setcounter{enumi}{2}
\item W.~A.~\textsc{Hiscock} and L.~\textsc{Lindblom},
  `Generic instabilities in first-order dissipative relativistic
  fluid theories',
  \emph{Phys.\ Rev.\ D}\ \textbf{31}, 725--33 (1985).

\item W.~A.~\textsc{Hiscock} and L.~\textsc{Lindblom},
  `Stability and causality in dissipative relativistic fluids',
  \emph{Ann.\ Phys.\ (N.Y.)}\ \textbf{151}, 466--96 (1983).

\item P.~\textsc{Romatschke},
  `New developments in relativistic fluid dynamics',
  \emph{Int.\ J.\ Mod.\ Phys.\ E}\ \textbf{19}, 1--53 (2010).
\end{enumerate}

\noindent The stability of relativistic viscous Taylor--Couette flow
(without axial flow) has been considered in:

\begin{enumerate}\setcounter{enumi}{5}
\item G.~S.~\textsc{Denicol}, T.~\textsc{Koide}, and
  D.~H.~\textsc{Rischke},
  `Dissipative relativistic fluid dynamics: a new way to derive
  the equations of motion from kinetic theory',
  \emph{Phys.\ Rev.\ Lett.}\ \textbf{105}, 162501 (2010).

\item M.~\textsc{Takamoto} and S.-i.~\textsc{Inutsuka},
  `The relativistic Kelvin--Helmholtz instability',
  \emph{J.\ Fluid Mech.}\ \textbf{735}, 347--65 (2013).
\end{enumerate}

\noindent The classical treatment of viscous spiral
Poiseuille--Couette flow follows:

\begin{enumerate}\setcounter{enumi}{7}
\item S.~\textsc{Chandrasekhar},
  `The hydrodynamic stability of viscid flow between coaxial
  cylinders',
  \emph{Proc.\ Nat.\ Acad.\ Sci.}\ \textbf{46}, 141--3 (1960).

\item R.~J.~\textsc{Donnelly} and \textsc{Dave Fultz},
  `Experiments on the stability of spiral flow between rotating
  cylinders',
  \emph{Proc.\ Nat.\ Acad.\ Sci.}\ \textbf{46}, 1150--4 (1960).
\end{enumerate}

\noindent For applications of causal viscous hydrodynamics
to astrophysical jets and accretion flows:

\begin{enumerate}\setcounter{enumi}{9}
\item J.~L.~\textsc{Zdunik}, P.~\textsc{Haensel}, and
  E.~\textsc{Gourgoulhon},
  `Recycled neutron stars with modified viscosity',
  \emph{Astron.\ Astrophys.}\ \textbf{372}, 535--43 (2001).

\item L.~\textsc{Rezzolla} and O.~\textsc{Zanotti},
  \emph{Relativistic Hydrodynamics},
  Oxford University Press, Oxford, 2013.

\item P.~\textsc{Romatschke} and U.~\textsc{Romatschke},
  \emph{Relativistic Fluid Dynamics In and Out of Equilibrium},
  Cambridge University Press, Cambridge, 2019.
\end{enumerate}
